\chapter{Пластичность, упрочнение и память}\label{ch:plasticity}

\section{Упругость, пластичность, разгрузка}

Диаграмма растяжения подробно разобрана в главе~\ref{ch:mech}. Здесь важно одно её свойство (рис.~\ref{fig:f7_ss}):
выше предела текучести $\sigma_y$ деформация уже не исчезает после снятия нагрузки. Разгрузка идёт по прямой с
тем же наклоном $E$, и остаётся пластическая деформация $\eps_p$. Чтобы деформировать металл дальше, нужно всё
большее напряжение --- это \textbf{упрочнение}. Полная деформация --- сумма упругой и пластической, $\eps = \sigma/E + \eps_p$,
и именно так считают пластичность и в модели гидридов.

\begin{figure}[htbp]\centering
\includegraphics[width=0.6\linewidth]{f7_ss}
\caption{Схема диаграммы растяжения: упругий участок (наклон $E$), предел текучести $\sigma_y$, упрочнение;
разгрузка упругая, остаётся пластическая деформация $\eps_p$.}
\label{fig:f7_ss}
\end{figure}

Для оболочечных сплавов циркония при 300--400\degC{} предел текучести порядка 200--500~МПа: меньше у
рекристаллизованного металла, больше у холоднодеформированного и снятого от напряжений.

\section{Когда металл течёт при сложном нагружении}

В трубе напряжения не одноосные: под давлением окружное $\sigma_\theta$ вдвое больше осевого $\sigma_z$. Нужен
критерий текучести для любого напряжённого состояния. Гидростатическое давление на текучесть почти не влияет ---
дислокации двигает только касательное напряжение. Поэтому критерий строят по девиатору.
\begin{keyf}[Критерий Мизеса]
\[
\sigma_{\text{экв}} = \sqrt{\tfrac32\,s_{ij}s_{ij}}\;=\;\sigma_y
\]
(для плоского состояния $\sigma_1^2-\sigma_1\sigma_2+\sigma_2^2 = \sigma_y^2$).
\end{keyf}
Критерий Треска --- максимальное касательное напряжение $(\sigma_{\max}-\sigma_{\min})/2 = \sigma_y/2$. На плоскости
главных напряжений Мизес --- эллипс, Треска --- вписанный шестиугольник (рис.~\ref{fig:f7_yield}).

\begin{figure}[htbp]\centering
\includegraphics[width=0.52\linewidth]{f7_yield}
\caption{Поверхности текучести в плоскости $(\sigma_\theta,\sigma_z)$: Мизес (эллипс) и Треска (шестиугольник).
Стрелки --- пути нагружения: одноосное растяжение, труба с заглушками ($\sigma_z = \sigma_\theta/2$), равное
двухосное. Труба начинает течь при $\sigma_\theta\approx1.155\,\sigma_y$.}
\label{fig:f7_yield}
\end{figure}

\begin{idea}
Почему при $\sigma_z = \sigma_\theta/2$ труба выдерживает больше, чем стержень? Осевое растяжение уменьшает
девиатор: два растяжения частично «гасят» друг друга по сдвигу. Поэтому одинаковые окружные напряжения в
трубе и в кольце, растягиваемом в одном направлении, дают разную пластичность --- и разную переориентацию.
\end{idea}

\section{Дислокационная природа текучести}

Дислокация начинает скользить, когда касательное напряжение в её плоскости и направлении скольжения достигает
критического $\tau_c$. При одноосном растяжении это напряжение равно $\sigma\cos\varphi\cos\lambda$ (закон
Шмида; $\varphi$ и $\lambda$ --- углы нагрузки с нормалью к плоскости и с направлением скольжения). Упрочнение
возникает потому, что дислокации множатся и мешают друг другу: $\tau\approx\alpha\mu b\sqrt\rho$, где $\rho$ ---
плотность дислокаций (уравнение Тейлора).

В цирконии лёгкое скольжение одно --- призматическое $\{10\bar10\}\langle11\bar20\rangle$. Деформация вдоль оси $c$
требует пирамидального скольжения $\langle c+a\rangle$ (оно в несколько раз «твёрже») или двойникования. Поэтому
прочность зерна сильно зависит от того, как к нагрузке повёрнута ось $c$. Прочность трубы в целом зависит от
текстуры (глава~\ref{ch:texture}). Анизотропную текучесть описывают критерием Хилла --- обобщением Мизеса с
разными весами для разных компонент напряжения.

\section{Скольжение в ГПУ: факторы Шмида и анизотропия}

Закон Шмида превращает таблицу систем скольжения (глава~\ref{ch:defects}) в карту: для каждого направления
растяжения можно найти наибольший фактор Шмида $m = \cos\varphi\cos\lambda$ среди систем данного типа
(рис.~\ref{fig:p1_schmid}). Призматическое скольжение хорошо ориентировано, когда нагрузка лежит в базисной
плоскости, и невозможно ($m = 0$) при нагрузке вдоль $c$. Базисное, наоборот, максимально при нагрузке под
45\textdegree{} к $c$. Пирамидальное $\langle c+a\rangle$ работает почти при любой ориентации, но его критическое
напряжение в разы выше.

\fig{p1_schmid}{Наибольший фактор Шмида при растяжении вдоль направления из стандартного треугольника ГПУ:
а) призматическое $\langle a\rangle$, б) базисное $\langle a\rangle$, в) пирамидальное $\langle c+a\rangle$ скольжение.
г) Распределение призматического фактора Шмида по зёрнам схематичной текстуры трубы (глава~\ref{ch:texture}) при
растяжении по окружности, по радиусу и по оси трубы.}

Для трубы это значит следующее (рис.~\ref{fig:p1_schmid}г). При окружном растяжении у большинства зёрен призматический
фактор Шмида 0.3--0.45 --- металл течёт легко. Деформация по радиусу (утонение стенки) требует скольжения в зёрнах,
у которых ось $c$ близка к радиусу; у них призматический фактор мал, и утонение стенки идёт трудно. Поэтому
оболочка с радиальной текстурой при растяжении сужается больше по ширине, чем по толщине (коэффициент Ланкфорда
$r = \eps_{\text{ширина}}/\eps_{\text{толщина}}$ больше единицы), а её поверхность текучести --- не эллипс Мизеса, а
вытянутый эллипс Хилла.

\textbf{Двойникование} дополняет скольжение там, где нужна деформация вдоль $c$. Двойник --- область, решётка которой
зеркально отражена относительно плоскости двойникования; при этом ось $c$ поворачивается почти на 85\textdegree{}
(двойники растяжения $\{10\bar12\}$). Двойникование включается при низких температурах и больших скоростях, и после него
локальная текстура резко меняется.

\section{Размер зерна: соотношение Холла--Петча}

Граница зерна мешает дислокации: в соседнем зерне другие системы скольжения, и дислокации скапливаются у границы.
Скопление концентрирует напряжение на границе тем сильнее, чем оно длиннее, то есть чем больше зерно. Отсюда
\begin{keyf}[Соотношение Холла--Петча]
\[
\sigma_y = \sigma_0 + k_y\,d^{-1/2}.
\]
\end{keyf}
Мелкое зерно упрочняет металл, не снижая пластичности, --- редкий случай. Для циркониевых сплавов с зерном в единицы
микрометров вклад границ составляет заметную часть предела текучести. Зерно меньше примерно 0.1~мкм (нанокристаллы)
уже не следует закону: дислокационные скопления не помещаются в зерно.

\section{Наклёп, возврат, рекристаллизация}

При холодной деформации плотность дислокаций растёт на порядки (глава~\ref{ch:defects}); металл упрочняется и
теряет пластичность (\textbf{наклёп}). В дислокациях запасается энергия --- около процента от работы деформации, но
именно она движет последующим отжигом. При нагреве идут три процесса (рис.~\ref{fig:p3_recryst}).
\begin{itemize}
\item \textbf{Возврат}: дислокации разных знаков аннигилируют, остальные выстраиваются в стенки (полигонизация).
  Внутренние напряжения снимаются, прочность снижается немного. Это и есть \textbf{отжиг для снятия напряжений} (для
  оболочек циркония --- около 470--500\degC): зёрна остаются вытянутыми, дислокаций много.
\item \textbf{Рекристаллизация}: в наклёпанном металле зарождаются и растут новые зёрна, почти свободные от
  дислокаций. Кинетика --- типичная кривая Аврами (глава~\ref{ch:nucleation}), $X = 1 - e^{-kt^n}$, с сильной
  зависимостью $k$ от температуры. Прочность падает до уровня отожжённого металла, пластичность восстанавливается.
  \textbf{Рекристаллизационный отжиг} циркониевых труб --- около 560--600\degC.
\item \textbf{Рост зерна}: после рекристаллизации крупные зёрна поедают мелкие, уменьшая площадь границ. Средний
  размер растёт как $d^2 - d_0^2 = kt$.
\end{itemize}

\fig{p3_recryst}{а) Доля рекристаллизованного объёма от времени отжига при трёх температурах (кинетика Аврами,
$n = 2$, энергия активации 300~кДж/моль; модельные параметры). б) Схема свойств после часового отжига
холоднодеформированного сплава: предел текучести, пластичность и доля рекристаллизации; СО --- отжиг для снятия
напряжений, РО --- рекристаллизационный.}

\begin{zrbox}
Состояние металла задаёт гидридам «рельеф». В холоднодеформированном и снятом от напряжений металле много
дислокаций и вытянутые зёрна: много мест зарождения, гидриды мелкие и густые, а межзёренные гидриды ориентированы
вдоль вытянутых границ. В рекристаллизованном металле зёрна равноосные, дислокаций мало, гидриды длиннее, и доля
межзёренных выше, а границы ориентированы во всех направлениях. Поэтому одно и то же напряжение даёт в этих
состояниях разную долю радиальных гидридов. В модели это входит через места зарождения, размер зерна и предел
текучести.
\end{zrbox}

\section{Пластическая аккомодация гидрида и память}

Упругие напряжения от гидрида --- порядка $\mu\eps^*\sim30~\text{ГПа}\cdot0.05\approx1.5$~ГПа. Это в несколько
раз выше предела текучести. Значит, вокруг всякого гидрида, кроме совсем мелкого, металл течёт. Оценим
размер пластической зоны для простейшей модели: жёсткий шар радиуса $a$ с линейным несоответствием $\delta$ в
идеально пластичной матрице (за зоной --- упругость, в зоне материал несжимаем). Объём, «выдавленный» шаром,
$4\pi a^3\delta$, должен равняться объёму, на который сдвигается упругая граница зоны радиуса $c$. Отсюда
\begin{keyf}[Пластическая зона вокруг включения]
\[
\Bigl(\frac{c}{a}\Bigr)^3 \approx \frac{6\mu\,\delta}{\sigma_y}.
\]
\end{keyf}
Для гидрида ($\mu = 33$~ГПа, $\delta\approx0.057$ --- треть объёмного несоответствия, $\sigma_y = 300$~МПа) зона в
3.4~раза больше самого гидрида (рис.~\ref{fig:f7_zone}). У настоящей пластинки зона вытянута вдоль неё и наиболее
сильна у торцов. Пластическое течение идёт рождением и выбросом дислокационных петель.

\fig{f7_zone}{а) Радиус пластической зоны вокруг несоразмерного включения в единицах его радиуса от предела
текучести матрицы при трёх несоответствиях. б) Схема: гидрид, пластический ореол и выброшенные дислокационные
петли.}

\begin{zrbox}
Пластическая аккомодация --- ключ к трём наблюдениям.
\begin{itemize}
\item \textbf{Гистерезис}: новый гидрид платит работой пластического течения, а при растворении эта работа не
  возвращается. Отсюда часть разрыва между TSSP и TSSD (глава~\ref{ch:solutions}).
\item \textbf{Память}: при растворении гидрида дислокации и остаточная деформация остаются. При следующем
  охлаждении гидрид выпадает там же --- места уже «подготовлены». Сколько памяти переживает выдержку при
  максимальной температуре, зависит от отжига (возврата) и проверяется опытами с многократными циклами.
\item \textbf{Коллективная пластичность}: когда гидридов много, их пластические зоны сливаются, и металл вокруг
  целой пачки пластинок течёт как одно целое. Это меняет поле напряжений, в котором зарождаются соседи.
\end{itemize}
\end{zrbox}

\section{Остаточные и межзёренные напряжения}

Даже без внешней нагрузки металл не свободен от напряжений. Их различают по масштабу:
\begin{itemize}
\item \textbf{I рода} --- на масштабе детали (после прокатки, сварки, изгиба);
\item \textbf{II рода} --- между зёрнами. В цирконии их создаёт анизотропия: при охлаждении зерно сжимается вдоль
  $c$ сильнее, чем поперёк, а соседи с другой ориентацией ему мешают. Оценка $\Delta\alpha\,\Delta T\,E$ даёт
  десятки МПа --- сопоставимо с нагрузками, при которых идёт переориентация;
\item \textbf{III рода} --- на масштабе дислокаций и включений (поле вокруг гидрида, глава~\ref{ch:elastic}).
\end{itemize}

\begin{zrbox}
Межзёренные напряжения разные в разных зёрнах, и каждое зерно «видит» внешнюю нагрузку со своей добавкой. Это
даёт разброс условий зарождения от зерна к зерну. Его величина --- один из параметров модели.
\end{zrbox}

\begin{exercise}
Найдите окружное напряжение начала текучести трубы с заглушками по Мизесу и по Треска. Насколько они различаются?
\end{exercise}

\begin{exercise}
Нагрузка приложена в базисной плоскости под углом $\psi$ к $\langle11\bar20\rangle$. Найдите наибольший призматический
фактор Шмида как функцию $\psi$. При каком $\psi$ он максимален?
\end{exercise}

\begin{exercise}
Выведите оценку размера пластической зоны $(c/a)^3 = 6\mu\delta/\sigma_y$. Для этого возьмите упругое решение
$u = B/r^2$ вне зоны, условие текучести на её границе и несжимаемость внутри.
\end{exercise}
